\chapter{作为有感行为统一理论的主动推断}
\section{引言}

本章将为主动推断收束全书的主要理论要点（第一部分）与实践实现（第二部分），并进一步把这些线索连接起来：我们将暂时抽离前面各章讨论的具体主动推断模型，转而聚焦这一框架的整合性面向。主动推断的一个重要优势在于，它为有感生物体必须解决的适应性问题提供了一套完整解法。因此，它能够以统一视角看待知觉、行动选择、注意、情绪调节等问题；而在心理学与神经科学中，这些问题通常被彼此分割，在人工智能中也往往由不同的计算方法分别处理。我们将把主动推断对这些问题的理解，放在控制论、动作观念运动理论、强化学习与最优控制等既有理论背景下加以讨论。最后，我们还会简要说明，主动推断的适用范围如何扩展到本书未深入展开的其他生物、社会与技术主题。

\section{总结前文}

本书系统介绍了主动推断的理论基础与实践实现。这里我们简要回顾前九章的讨论，以便重温后续分析所需的关键概念。

在第 1 章中，我们将主动推断引入为一种规范性方法，用来理解那些与环境构成行动--知觉回路的有感生物体（Fuster 2004）。所谓规范性方法，是从第一性原理出发，推导并检验有关目标现象的经验预测；在这里，目标现象就是生物体如何在与环境进行适应性交换（行动--知觉回路）的同时持续存在。我们还指出，人们可以沿着一条“低路”或一条“高路”抵达主动推断。

在第 2 章中，我们展示了通往主动推断的低路。这一路径始于这样一种观念：大脑是一台预测机器，拥有一个生成模型，即关于世界中隐变量如何生成感觉输入的概率性表征（例如，苹果反射的光如何刺激视网膜）。通过对这一模型做反演，大脑便能推断感觉的成因（例如，在视网膜以某种方式受到刺激的情况下，我是否看到了一个苹果）。这种关于知觉的看法，也即“知觉即推断”，其历史根源可以追溯到 Helmholtz 的无意识推断观，以及更晚近的贝叶斯大脑假说。主动推断进一步把动作控制与规划也纳入推断的范围，即“控制即推断”“规划即推断”。更重要的是，它表明知觉与行动并不是本质上彼此分离的过程，而是在实现同一个目标。我们首先以较为非正式的方式描述这一目标，把它说成是最小化模型与世界之间的不一致，而这通常可归结为最小化惊异或预测误差。简单说来，最小化模型与世界之间的不一致有两种方式：要么改变自己的心智以适配世界（知觉），要么改变世界以适配模型（行动）。这些都可以用贝叶斯推断来描述。然而，精确推断往往难以直接求解，因此主动推断采用了一种（变分）近似，并注意到精确推断本身也可被视为近似推断的特例。由此，我们得到了对知觉与行动共同目标的第二种、也更正式的描述：最小化变分自由能。这是主动推断的核心量，可以被拆解为若干组成部分，例如能量与熵、复杂性与准确性，或惊异与散度。最后，我们又引入了第二种自由能：期望自由能。它在规划过程中尤其重要，因为它为替代策略打分提供了方法，即考虑这些策略预期会生成的未来结果。它同样可以展开为若干组成部分，例如信息增益与务实价值，或期望模糊性与风险。

在第 3 章中，我们展示了通往主动推断的高路。这条替代路径从一个“去通胀”的生物学要求出发：生物体必须维持自身完整性并避免耗散，而这可以描述为避免落入令人惊异的状态。随后我们引入了马尔可夫毯的概念，它形式化了生物体内部状态与世界外部状态之间的统计分离。关键在于，内部状态与外部状态只能通过中间变量，也就是毯状态（主动状态与感觉状态）间接影响彼此。这种由马尔可夫毯介导的统计分离，对于赋予生物体相对于外部世界的一定自主性至关重要。要理解这种视角为何有用，可以看它带来的三个后果。第一，具有马尔可夫毯的生物体，在贝叶斯意义上仿佛拥有一个关于外部环境的模型：它的内部状态，平均而言，对应于关于世界外部状态的近似后验信念。第二，这种自主性之所以能够成立，是因为该生物体的模型（即其内部状态）并非无偏，而是预设了一些必须维持的生存前提，或者说先验偏好；例如，对鱼来说，就是待在水里。第三，在这一形式之下，我们可以把相对于先验偏好的最优行为描述为：通过知觉与行动最大化（贝叶斯）模型证据。通过最大化模型证据，也即“自证”，生物体确保自己实现先验偏好（例如鱼留在水中），并避开令人惊异的状态。反过来，最大化模型证据在数学上与最小化变分自由能近似等价，因此我们再次以另一条路径抵达了第 2 章中的核心构造。最后，我们还详细讨论了最小化惊异与 Hamilton 最小作用原理之间的关系，这揭示了主动推断与统计物理第一性原理之间的形式联系。

在第 4 章中，我们概述了主动推断的形式方面。重点是从贝叶斯推断走向可计算的近似，也就是变分推断，以及生物体由此通过知觉与行动最小化变分自由能的目标。这里最重要的启发是：生物体用来理解世界的生成模型至关重要。我们介绍了两类生成模型，它们表达我们关于数据如何被生成的信念，分别使用离散变量或连续变量。两者都支持主动推断，但分别适用于以离散时间形式表述的事态（如部分可观测马尔可夫决策问题）与以连续时间形式表述的事态（如随机微分方程）。

在第 5 章中，我们强调了作为规范性原理的自由能最小化，与关于大脑如何实现这一原理的过程理论之间的区别，并说明后者会生成可检验的预测。随后我们概述了伴随主动推断而来的若干过程理论，它们涵盖神经元消息传递、神经解剖回路（如皮层--皮层下回路）以及神经调质等领域。例如，在解剖层面，消息传递与经典皮层微回路高度契合：在某一级，来自深层皮层的预测会投向下一层级的浅层皮层（Bastos et al. 2012）。在更系统的层面，我们讨论了贝叶斯推断、学习与精度加权分别如何对应神经动力学、突触可塑性与神经调质，以及预测编码中的自上而下与自下而上的神经消息传递如何映射到较慢的脑节律（如 alpha 或 beta）与较快的脑节律（如 gamma）。这些以及其他例子表明，一旦设计出一个特定的主动推断模型，就可以从其生成模型的形式出发，导出神经生物学层面的含义。

在第 6 章中，我们给出了一套设计主动推断模型的配方。我们看到，尽管所有生物体都在最小化自身的变分自由能，但由于它们拥有不同的生成模型，所以会表现出彼此不同、甚至相反的行为。因此，区分不同生物体（例如较简单者与较复杂者）的关键，就在于它们的生成模型。可能的生成模型有着丰富谱系，它们对应不同的生物学实现（例如神经实现），并会在不同情境与生态位中产生不同的适应性或失适应性行为。正因如此，主动推断既适合刻画那种能够感知并追逐营养梯度的简单生物（如细菌），也适合刻画像我们这样追求复杂目标并参与丰富文化实践的复杂生物，甚至适用于不同个体，只要能够恰当地刻画各自的生成模型即可。进化似乎逐步发现了越来越复杂的脑和身体设计结构，使生物体能够应对并塑造更丰富的生态位。建模者则可以逆向工程这一过程，依据目标生物所占据的生态位，以生成模型的形式指定其脑与身体的设计。这对应于一系列设计选择，例如使用离散变量还是类别变量，使用浅层模型还是层级模型；这些我们都在该章中进行了展开。

在第 7 章与第 8 章中，我们提供了大量离散时间与连续时间下的主动推断模型示例，涉及知觉推断、目标导向导航、模型学习、动作控制等问题。这些例子意在展示这类模型下涌现行为的多样性，并详细说明它们在实践中如何被具体设定。

在第 9 章中，我们讨论了如何使用主动推断进行基于模型的数据分析，以及如何恢复一个个体生成模型中最能解释其任务行为的参数。这种计算表型分析使用的仍然是本书其余部分讨论过的那种贝叶斯推断，只是用途不同：它帮助我们设计并评估关于他人主观模型的客观模型。

\section{连接各个线头：主动推断的整合性视角}

数十年前，哲学家 Dennett 曾抱怨，认知科学家把太多精力花在对孤立子系统建模上，例如知觉、语言理解等，而这些子系统的边界往往带有很大任意性。他建议，与其如此，不如试着去建模“整只鬣蜥”：一个完整的认知生物体，也许相对简单，但要连同它所处、必须应对的环境生态位一起建模（Dennett 1978）。

主动推断的一个重要优点，就是它从第一性原理出发解释生物体如何解决适应性问题。本书所采取的规范性路线假定，我们可以从变分自由能最小化这一原理出发，推导出关于具体认知过程及其神经基础的含义，例如知觉、行动选择、注意与情绪调节。

设想一个简单生物，它必须解决寻找食物或庇护所之类的问题。若以主动推断来表述，这类问题可以用具身行动的语言描述为：通过行动去召唤偏好的感觉，例如与食物相关的感觉。只要这些偏好感觉已作为先验信念写入其生成模型，那么这个生物体实际上就在为自己的模型收集证据，或者更形象地说，为自己的存在收集证据，也就是最大化模型证据，进行自证。这个简单原理会延伸到通常被分开处理的多种心理功能，包括知觉、动作控制、记忆、注意、意向、情绪等。比如，知觉与行动都属于自证行为：生物体既可以通过改变信念来让自身期望与感觉对齐，例如更新关于“是否有食物”的信念；也可以通过改变世界来做到这一点，例如主动获得与食物相关的感觉。记忆与注意同样可以被看作优化同一目标。长期记忆通过学习生成模型的参数而发展；工作记忆则是信念更新，只不过这些信念关涉外部状态的过去与未来；注意是对感觉输入精度信念的优化。各种形式的规划与意向性，则可以通过某些生物体在若干可能未来之间进行选择的能力来理解，而这又要求其拥有时间深度生成模型。这类模型能够预测某条行动路线会产生怎样的结果，并且对这些结果持一种乐观态度。这种乐观体现为：相信未来结果将导向偏好结果。时间深度模型还帮助我们理解更复杂的前瞻与回顾形式：前者是以关于现在的信念推导关于未来的信念，后者则是以关于现在的信念更新关于过去的信念。各种内感受调节与情绪，也可以诉诸关于内部生理状态的生成模型来理解；这类模型会预测未来事件的异稳态后果。

上述例子表明，从一种关于有感行为的规范性理论来研究认知与行为，会带来一个重要后果。这样的理论并不是先把知觉、决策、规划等分离的认知功能拼装起来；相反，它先为生物体必须解决的问题提供一个完整解法，然后再分析这一解法，从中导出关于各类认知功能的含义。譬如，究竟是什么机制使一个生命体或人工体（如机器人）能够知觉世界、记住它，或对其进行规划（Verschure et al. 2003, 2014; Verschure 2012; Pezzulo, Barsalou et al. 2013; Krakauer et al. 2017）？这一转向很重要，因为心理学与神经科学教科书中常用的认知功能分类，大多继承自早期哲学与心理学理论，有时也被称为 James 式范畴。它们虽然具有很高的启发价值，却可能相当任意，或者根本不对应彼此分离的认知与神经过程（Pezzulo and Cisek 2016; Buzsaki 2019; Cisek 2019）。的确，这些 James 式范畴或许更像是我们用来解释自己与感觉世界互动的生成模型成分，而不是真正解释这种互动本身的机制。举例来说，那种带有唯我论色彩的假设“我正在知觉”，也许只是我对当前事态的一种解释，而这些当前事态中就包括了我正在更新信念这一事实。

采取规范性视角还有助于识别不同领域中认知现象之间的形式类比。一个例子就是“探索与利用的权衡”，它以多种形式反复出现（Hills et al. 2015）。这一权衡常常在觅食情境中被研究：生物体必须在利用以往成功方案与探索新的、可能更优的方案之间做出选择。然而，同样的权衡也会在记忆搜索与资源受限的审议过程中出现，例如当时间有限或搜索努力受限时，生物体必须在继续利用当前最佳方案，与投入更多时间和认知努力以探索更多可能性之间抉择。若用自由能来刻画这些看似无关的现象，就可能揭示其深层相似性（Friston, Rigoli et al. 2015; Pezzulo, Cartoni et al. 2016; Gottwald and Braun 2020）。

最后，除了为心理现象提供统一视角，主动推断还为理解其对应的神经计算提供了原则化方法。换言之，它提供了一种过程理论，把认知处理与预期的神经动力学联系起来。主动推断假定，凡是与大脑、心智和行为相关的重要内容，都可用变分自由能最小化来描述。反过来，这种最小化会带来特定的神经签名，例如消息传递样式或脑解剖结构，而这些都可以接受经验检验。

本章余下部分将像是在勾勒一本心理学教材那样，讨论主动推断对若干心理功能的含义。针对每一种功能，我们也会指出主动推断与文献中其他流行理论之间的联系与分歧。

\section{预测性大脑、预测性心智与预测加工}

传统的大脑与认知理论强调从外部刺激到内部表征、再到运动行动的前馈式转换。这种观点被称为“三明治模型”，因为凡是夹在刺激与反应之间的内容都被统一贴上“认知”的标签（Hurley 2008）。在这一视角下，大脑的主要功能是把输入刺激转化为与情境相适配的反应。

主动推断与此显著不同，它强调大脑与认知的预测性与目标导向性。用心理学语言说，主动推断生物体（或其大脑）是一种概率推断机器，会基于自身生成模型持续地产生预测。

进行自证的生物体以两种基本方式使用这些预测。第一，它们把预测与到来的数据进行比较，以检验自身假设（预测编码），并在更慢的时间尺度上修正模型（学习）。第二，它们会把预测付诸实施，以引导自身如何去采样数据（主动推断）。通过这样做，主动推断生物体履行了两个要求：其一是认识性的，例如去视觉探索那些存在显著信息、能够消解关于假设或模型不确定性的地方；其二是务实性的，例如移动到能够确保偏好观测、如奖励的位置。认识性要求使知觉与学习都成为主动过程；务实性要求则使行为具有目标导向性。

\subsection{预测加工}

这种以预测与目标为中心的大脑与认知观，与预测加工（predictive processing, PP）密切相关，也为其提供了重要灵感。预测加工是心灵哲学与认识论中的一个新兴框架，强调预测在大脑与认知中的核心地位，并诉诸“预测性大脑”或“预测性心智”等概念（Clark 2013, 2015; Hohwy 2013）。

预测加工理论有时会诉诸主动推断的具体机制及其若干核心构件，例如生成模型、预测编码、自由能、精度控制与马尔可夫毯；但有时它们也会诉诸其他并不属于主动推断的构件，例如成对耦合的逆模型与前向模型。因此，“预测加工”这一术语的使用范围比“主动推断”更广，也更少受约束。

预测加工理论在哲学中引起了相当大的关注，因为它在多重意义上都具有统一潜力：它可以横跨知觉、行动、学习与精神病理等多个认知领域；可以贯通从较低层级的感知运动处理到较高层级的心理构造；还可以从简单生物体一路延伸到大脑、个体以及社会和文化构造。预测加工理论的另一种吸引力在于，它使用了信念、惊异等概念性术语，这些术语与哲学家熟悉的心理学分析层面相衔接。当然，需要注意的是，这些术语有时带有不同于日常用法的技术含义。

然而，随着预测加工受到越来越多关注，人们也逐渐发现，哲学家们对于它的理论与认识论含义并不一致。例如，有人将其解读为内在主义的（Hohwy 2013），有人强调其具身性或行动基础（Clark 2015），也有人以生成论或非表征主义的方式来理解它（Bruineberg et al. 2016; Ramstead et al. 2019）。围绕这些概念解释的争论超出了本书讨论范围。

\section{知觉}

主动推断把知觉视为一种基于生成模型的推断过程，这个生成模型描述感觉观测是如何产生的。贝叶斯规则本质上是在对这一模型进行反演，以在给定观测的情况下计算关于环境隐状态的信念。把知觉理解为推断的思想可追溯到 Helmholtz（1866），并在心理学、计算神经科学与机器学习中多次被重新提出，例如“分析即综合”等观点（Gregory 1980; Dayan et al. 1995; Mesulam 1998; Yuille and Kersten 2006）。这种生成建模取向已被证明能够有效应对困难的知觉问题，例如破解基于文本的 CAPTCHA（George et al. 2017）。

\subsection{贝叶斯大脑假说}

这一思想在当代最突出的表达，就是贝叶斯大脑假说，它已被应用于决策、感觉处理与学习等多个领域（Doya 2007）。主动推断为这些推断性思想提供了规范性基础，因为它把它们从变分自由能最小化这一要求中推导出来。由于同一要求也扩展到行动动力学，主动推断自然能够建模主动知觉，以及生物体如何通过主动采样观测来检验自己的假设（Gregory 1980）。相比之下，在贝叶斯大脑的议程下，知觉与行动通常依照不同的要求来建模，其中行动需要诉诸贝叶斯决策理论，见第 \ref{sec:bayesian-decision-theory} 节。

更广义地说，贝叶斯大脑假说指的是一族彼此未必整合、且常常给出不同经验预测的方法。例如，它包括计算层级的主张，即大脑执行贝叶斯最优的感觉运动与多感觉整合（Kording and Wolpert 2006）；算法层级的主张，即大脑实现某种特定的贝叶斯近似，例如抽样式决策（Stewart et al. 2006）；以及神经层级的主张，即神经群体如何执行概率计算或编码概率分布，例如用样本或概率性群体编码来实现（Fiser et al. 2010; Pouget et al. 2013）。在每一个解释层级上，都存在相互竞争的理论。例如，人们常诉诸精确贝叶斯推断的近似来解释偏离最优行为的现象，但不同工作考虑的是不同的、且并不总彼此兼容的近似方案，例如不同的抽样方法。更广泛地说，不同层级上的提案之间也并不总能顺畅对接。这是因为，贝叶斯计算可以通过多种算法路径来实现或近似，即便系统并未显式表征概率分布（Aitchison and Lengyel 2017）。

主动推断提供了一个更加整合的视角，它把规范性原理与过程理论连接起来。在规范性层面，其核心假设是所有过程都在最小化变分自由能。与之相应的推断过程理论，是对自由能做梯度下降；这一点具有明确的神经生理学含义，我们已在第 5 章中讨论过（Friston, FitzGerald et al. 2016）。更广泛地说，人们可以从自由能最小化原理出发，推导关于脑结构的含义。

例如，知觉推断在连续时间中的经典过程模型就是预测编码。预测编码最初由 Rao 与 Ballard（1999）提出，作为一种层级化知觉处理理论，用以解释一系列已知的自上而下效应；这些效应难以与纯前馈架构以及既有生理事实（例如感觉层级中前向/自下而上与后向/自上而下连接的存在）相调和。然而，在某些假设下，如 Laplace 近似，预测编码可以从自由能最小化原理中推导出来（Friston 2005）。进一步地，连续时间中的主动推断可以被构造为对预测编码向行动领域的定向扩展，即给一个预测编码主体配上运动反射（Shipp et al. 2013）。这就引出了下一节。

\section{动作控制}

在主动推断中，行动处理与知觉处理是平行的，因为两者都由前向预测引导，只不过前者主要依赖本体感觉预测，后者主要依赖外感受预测。正是“我的手抓住了杯子”这一类本体感觉预测，诱发了抓取动作。行动与知觉的等价性在神经生物学层面同样成立：运动皮层的组织方式与感觉皮层相同，也是一种预测编码架构；不同之处在于，它能够影响脑干与脊髓中的运动反射（Shipp et al. 2013），而且它接收的上行输入相对较少。运动反射使系统可以通过设定沿着期望轨迹的一系列“平衡点”来控制运动，这一思想与平衡点假说相呼应（Feldman 2009）。

需要强调的是，启动一个动作，例如抓杯子，要求对先验信念与感觉流的精度（方差的倒数）进行适当调节。这是因为，这些精度的相对大小决定了生物体如何处理以下冲突：一方面，它的先验信念认为自己握住了杯子；另一方面，感觉输入却表明它并未握住。若关于“抓住杯子”的先验信念精度很低，那么它就很容易在相冲突的感觉证据面前被修正，结果是改变想法，而不发生动作。相反，当先验信念占据主导地位，也就是具有更高精度时，它即便面对冲突性感觉证据仍会被维持下来，从而诱发抓取动作，以解决这一冲突。为了确保这一点，动作启动会引发短暂的感觉衰减，也就是下调感觉预测误差的权重。如果这种感觉衰减失效，便可能带来失适应后果，例如无法启动或控制运动（Brown et al. 2013）。

\subsection{动作观念运动理论}

在主动推断中，动作源自本体感觉预测，而不是运动指令（Adams, Shipp, and Friston 2013）。这一思想把主动推断与动作观念运动理论联系起来。该理论可追溯到 William James（1890），后来又发展出“事件编码”与“预期行为控制”等理论（Hommel et al. 2001; Hoffmann 2003）。动作观念运动理论认为，动作--效果联结（类似前向模型）是认知架构中的关键机制。重要的是，这些联结可以双向使用：当它们沿着“动作到效果”的方向使用时，可以生成感觉预测；当它们沿着“效果到动作”的方向使用时，则可以选择那些能实现预期知觉后果的动作。这意味着，动作的选择与控制，是基于它们预期后果来进行的，这也正是“观念 + 运动”这一名称的含义。这种前瞻性的动作控制观得到了大量文献支持，这些研究记录了预期动作后果对动作选择与执行的影响（Kunde et al. 2004）。主动推断为这一思想提供了数学刻画，并补入了若干额外机制，例如精度控制与感觉衰减的重要性；这些机制在动作观念运动理论中尚未被充分展开，但与其是兼容的。

\subsection{控制论}

主动推断与控制论关于行为具有目的性、目标导向性以及主体--环境反馈交互重要性的观点密切相关，这一点在 TOTE（Test, Operate, Test, Exit）及其相关模型中表现得尤为典型（Miller et al. 1960; Pezzulo, Baldassarre et al. 2006）。在 TOTE 与主动推断中，动作的选择都由偏好目标状态与当前状态之间的差异所决定。这些方法都不同于简单的刺激--反应关系，后者更常见于行为主义理论与强化学习等计算框架中（Sutton and Barto 1998）。

主动推断中的动作控制观，尤其接近知觉控制理论（perceptual control theory）（Powers 1973）。该理论的核心主张是：被控制的是知觉状态，而不是运动输出或动作本身。比如开车时，我们所控制并在扰动下努力维持稳定的，是一个参考速度或目标速度，例如时速 90 英里，这一目标通过速度表反映出来；至于为了保持这一速度而采取的动作，例如加速还是减速，则更具可变性，也更依赖情境。例如，根据扰动的性质不同，如风、陡坡或其他车辆，我们可能需要加速，也可能需要减速，才能维持参考速度。这一看法实现了 William James（1890）的一个建议：人类是以灵活手段达成稳定目标的。

尽管在主动推断与知觉控制理论中，真正控制动作的都是知觉预测，尤其是本体感觉预测，但两者在控制如何实施这一点上有所不同。主动推断而非知觉控制理论，强调动作控制具有基于生成模型的前瞻性或前馈成分。相比之下，知觉控制理论认为反馈机制基本已足够控制行为，而试图预测扰动、实施前馈控制或开环控制是没有价值的。不过，这一反对意见主要针对的是那些使用逆--前向模型的控制理论的局限，见下一小节。在主动推断中，生成模型或前向模型并不是用来预测扰动，而是用来预测需要通过行动来实现的未来期望状态与轨迹，并据此推断知觉事件的潜在原因。

主动推断与知觉控制理论之间另一重要契合点，是它们对控制层级的理解。知觉控制理论提出，更高层级通过设定更低层级的参考点或设定点，也就是规定低层级“要实现什么”，来控制低层级；而低层级仍可自由选择实现目标的手段，而不是被直接指定或偏置“怎么做”。这与大多数关于层级控制与自上而下控制的理论形成对照，后者通常认为高层级要么直接选择计划（Botvinick 2008），要么偏置低层级对动作或运动指令的选择（Miller and Cohen 2001）。类似知觉控制理论，主动推断也允许把层级控制分解为一个自上而下的目标与子目标级联结构，并使之在适当的较低层级上被自主实现。此外，在主动推断中，控制层级不同层次上表征的目标，其贡献还可以被动机过程所调制，即按精度加权，从而优先处理更显著或更紧迫的目标（Pezzulo, Rigoli, and Friston 2015, 2018）。

\subsection{最优控制理论}

主动推断对动作控制的说明，与神经科学中其他控制模型，如最优控制理论，有显著差异（Todorov 2004; Shadmehr et al. 2010）。该框架认为，大脑的运动皮层通过一个把刺激映射到反应的（反应式）控制策略来选择动作。而主动推断则认为，运动皮层传递的是预测，而不是指令。

此外，虽然最优控制理论与主动推断都诉诸内部模型，但它们对内部建模的理解不同（Friston 2011）。在最优控制中，内部模型分为两类：逆模型编码刺激--反应联结，并依据某种代价函数选择运动指令；前向模型则编码动作--结果联结，并向逆模型提供模拟输入，以替代噪声较大或延迟的反馈，从而超越纯反馈控制。逆模型与前向模型还可以在一个脱离外部行动--知觉回路的闭环中运作，也就是当输入与输出被抑制时，为内部的“如果如此会怎样”的动作序列模拟提供支持。这种内部动作模拟被用于解释多种认知功能，例如规划、动作知觉与社会领域中的模仿（Jeannerod 2001; Wolpert et al. 2003），以及多种运动障碍与精神病理现象（Frith et al. 2000）。

与前向--逆向建模方案不同，在主动推断中，动作控制的主要负担由前向（生成）模型承担，而逆模型非常简约，通常可以简化为在外周层级，也就是脑干或脊髓中解决的简单反射。当预期状态与观测状态之间存在差异时，例如期望的手臂位置与当前手臂位置不同，动作便被启动；这也就是一种感觉预测误差。换言之，在主动推断中，运动指令等价于前向模型作出的预测，而不是像最优控制中那样由逆模型计算出的结果。本体感觉预测误差会通过动作，也就是手臂移动，被消解。主动推断认为，这一需要由动作来填补的缺口很小，因此无需复杂的逆模型，只需简单得多的运动反射即可（Adams, Shipp, and Friston 2013）。\footnote{运动反射之所以比逆模型更简单，是因为它不编码从推断出的世界状态到动作的映射，而只编码动作与感觉后果之间更简单的映射。关于这一点的进一步讨论，见 Friston, Daunizeau et al. (2010)。}

最优运动控制与主动推断之间另一个关键差异在于，前者使用代价函数或价值函数来驱动行动，而后者则以贝叶斯意义上的先验来取而代之，更准确地说，是以期望自由能中隐含的先验偏好来承担这一角色。下一节将继续讨论这一点。

\section{效用与决策}

状态代价函数或价值函数的观念，在最优运动控制、效用最大化的经济学理论与强化学习等许多领域都居于核心地位。例如，在最优控制理论中，针对一个到达任务的最优控制策略，常被定义为使某一特定代价函数最小的策略，例如运动最平滑或最小跃度。在强化学习问题中，例如在带有一个或多个奖励的迷宫中导航，最优策略则是那个既能最大化贴现奖励，又能最小化运动成本的策略。这些问题往往通过 Bellman 方程来求解，在连续时间中则是 Hamilton--Jacobi--Bellman 方程。其一般思想是：一个决策的价值可以分解成两部分，即即时奖励与剩余决策问题的价值。正是这种分解使动态规划这一迭代程序成为可能，而动态规划正是控制理论与强化学习的核心（Bellman 1954）。

主动推断与上述思路在两个主要方面不同。第一，主动推断并不只考虑效用最大化，而是考虑更广义的目标，即最小化（期望）自由能；其中还包含额外的认识性要求，例如消除当前状态歧义与寻求新奇，见图 2.5。这些额外目标有时会在经典奖励之外被额外加入，例如作为“新奇奖励加成”（Kakade and Dayan 2002）或“内在奖励”（Schmidhuber 1991; Oudeyer et al. 2007; Baldassarre and Mirolli 2013; Gottlieb et al. 2013）；但在主动推断中，它们会自动出现，使系统能够处理许多决策中隐含的探索--利用权衡。原因在于，自由能是关于信念的泛函，因此我们进入的是信念优化而非外部奖励函数的领域；这一点在探索性问题中尤其关键，因为在这类问题里，成功依赖于尽可能消除不确定性。

第二，在主动推断中，代价的观念被吸收到先验之中。先验，或先验偏好，规定了控制的目标，例如要遵循的轨迹或要抵达的终点。用先验来编码偏好观测或偏好观测序列，可能比使用效用更具表达力（Friston, Daunizeau, and Kiebel 2009）。在这种方法下，寻找最优策略被改写成一个推断问题：推断出那组能够实现偏好轨迹的控制状态序列。它不需要显式价值函数或 Bellman 方程，尽管可以诉诸类似的递归逻辑（Friston, Da Costa et al. 2020）。至少有两个基本差异，区分了主动推断与强化学习中通常使用先验与价值函数的方式。其一，强化学习使用的是状态价值函数或状态--动作对价值函数，而主动推断使用的是关于观测的先验。其二，价值函数是以遵循特定策略时，处于某个状态（或在某个状态下采取某一动作）的期望回报来定义的，即从该状态出发并执行该策略后所获得的未来贴现奖励之和。相比之下，在主动推断中，先验通常既不把未来奖励相加，也不对其贴现。相反，只有在计算某一策略的期望自由能时，才会涌现出与期望回报相近的量。因此，期望自由能才是价值函数最接近的对应物。但即便如此，它也不同，因为期望自由能是关于状态信念的泛函，而不是状态本身的函数。尽管如此，人们仍然可以构造出近似于强化学习状态价值函数的先验，例如把期望自由能的计算结果缓存到这些状态之中（Friston, FitzGerald et al. 2016; Maisto, Friston, and Pezzulo 2019）。

此外，把效用吸收到先验中还带来一个重要的理论后果：先验承担了目标的角色，并使生成模型具有偏置，或者说具有一种乐观性，即生物体相信自己会遭遇偏好的结果。正是这种乐观性支撑了主动推断中那些能实现期望结果的被推断计划；若缺乏这种类型的乐观，便可能对应于冷漠或意志缺失（Hezemans et al. 2020）。这与贝叶斯决策理论等其他决策形式化方法形成对照，因为那些方法会把事件的概率与其效用区分开来。不过，这种区分在某种程度上只是表面的，因为一个效用函数总是可以被改写为某种先验信念，这与如下事实一致：最大化该效用函数的行为，先天地、按设计就是更可能发生的。从一种稍显同义反复式的去通胀视角看，这本就是“效用”的定义。

\subsection{贝叶斯决策理论}\label{sec:bayesian-decision-theory}

贝叶斯决策理论是一套数学框架，它把贝叶斯大脑的思想扩展到决策、感觉运动控制与学习等领域（Kording and Wolpert 2006; Shadmehr et al. 2010; Wolpert and Landy 2012）。它将决策描述为两个相互区分的过程。第一过程利用贝叶斯计算预测未来结果的概率，这些结果依赖于动作或策略；第二过程则利用固定或学习到的效用函数或代价函数，来定义对计划的偏好。最终的决策过程，也即动作选择，会整合这两条信息流，从而以更高概率选择那条更可能带来更高回报的行动方案。这与主动推断形成对比：在主动推断中，先验分布直接标示出对生物体而言什么是有价值的，或者说什么在进化历史上曾经有价值。不过，人们仍可在贝叶斯决策理论的这两条信息流，与主动推断中变分自由能与期望自由能的优化之间建立类比。在主动推断里，最小化变分自由能能够带来关于世界状态及其可能演化的准确且简洁的信念；而通过策略选择来最小化期望自由能的先验信念，则内含了偏好的观念。

在某些学术圈中，人们对贝叶斯决策理论的地位有所疑虑。这一疑虑源自完全类定理（Wald 1947; Brown 1981）：对于任意给定的一组决策与代价函数，总存在某些先验信念使这些决策成为贝叶斯最优。也就是说，当先验信念与代价函数被分开处理时，两者之间存在一种隐含的对偶性或退化性。在某种意义上，主动推断正是通过将效用或代价函数吸收进先验信念并表述为偏好，来消解这种退化。

\subsection{强化学习}

强化学习（reinforcement learning, RL）是一种求解马尔可夫决策问题的方法，在人工智能与认知科学中都极为流行（Sutton and Barto 1998）。它关注主体如何通过试错来学习一个策略，例如平衡杆子的策略；主体尝试动作，例如向左移动，并根据动作成功（如杆子保持平衡）或失败（如杆子倒下）而获得正或负的强化。

主动推断与强化学习处理的问题集有重叠，但在数学形式与概念理解上有许多不同。正如前文所述，主动推断抛弃了奖励、价值函数与 Bellman 最优性这些强化学习的核心概念。此外，“策略”这一概念在两种框架中的用法也不同。在强化学习中，策略表示一组需要通过学习获得的刺激--反应映射；而在主动推断中，策略是生成模型的一部分，表示需要被推断出来的控制状态序列。

强化学习方法种类繁多，但大致可分为三大家族。前两类方法都试图学习良好的价值函数，只是路径不同。

无模型强化学习方法直接从经验中学习价值函数：它们执行动作、收集奖励、更新价值函数，再利用价值函数更新策略。之所以称为“无模型”，是因为它们不使用那种能够预测未来状态的转移模型，而这种模型正是主动推断所使用的。取而代之，它们隐含地诉诸更简单的模型，例如状态--动作映射。无模型强化学习中的价值函数学习，往往涉及奖励预测误差的计算，例如广为人知的时序差分规则。主动推断当然也常诉诸预测误差，但那是状态预测误差，而非奖励预测误差，因为主动推断里没有奖励这一概念。

有模型强化学习方法并不直接从经验中学习价值函数或策略。相反，它们从经验中学习一个任务模型，利用该模型进行规划，也就是模拟可能经验，然后再用这些模拟经验来更新价值函数与策略。虽然主动推断与强化学习都诉诸基于模型的规划，但二者使用方式不同。在主动推断中，规划是为每一条策略计算期望自由能，而不是用来更新价值函数。严格说来，如果把期望自由能看作一种价值泛函，那么也可以说，借助生成模型得出的推断结果被用于更新这个泛函；这为两者提供了一个类比点。

第三类强化学习方法是策略梯度法。它们试图直接优化策略，而不经过价值函数这一中间层，而价值函数恰恰是有模型与无模型强化学习的核心。这类方法从参数化策略出发，这些策略能够生成例如运动轨迹，然后通过调整参数来优化策略：如果某条轨迹带来较高的正奖励，就提高该策略发生的概率；反之，则降低其概率。这一点使策略梯度方法与主动推断存在某种关联，因为主动推断同样不依赖价值函数（Millidge 2019）。但策略梯度的一般目标，即最大化长期累积奖励，与主动推断的目标仍然不同。

除了形式差异之外，主动推断与强化学习在概念上也有若干重要区别。其中一个区别涉及对目标导向行为与习惯行为的解释。在动物学习文献中，目标导向选择由关于动作与结果之间联结的前瞻性知识所中介（Dickinson and Balleine 1990）；而习惯性选择则不具前瞻性，依赖更简单的机制，例如刺激--反应机制。强化学习中的一个流行观点认为，目标导向与习惯性选择分别对应于有模型与无模型强化学习，并且二者是并行习得、持续竞争行为控制权的（Daw et al. 2005）。

主动推断则把目标导向与习惯性选择映射到不同机制。在离散时间的主动推断中，策略选择本质上是基于模型的，因此符合目标导向、审议性选择的定义。这一点与有模型强化学习相似，但仍有区别：在有模型强化学习中，动作是在前瞻意义上被选择的，因为使用了模型，但动作控制本身仍是反应式的，依赖刺激--反应策略；而在主动推断中，动作可以通过实现本体感觉预测而以前摄方式受到控制，参见第 10.6 节。

在主动推断中，习惯可以通过反复执行目标导向策略，并缓存“在什么情境下哪些策略会成功”的信息来获得。这些缓存的信息可以作为策略的先验价值被并入模型（Friston, FitzGerald et al. 2016; Maisto, Friston, and Pezzulo 2019）。这一机制允许系统在不经审议的情况下，直接执行在给定情境中具有高先验价值的策略。简单说，这相当于主体反复观察“我在做什么”，进而学到“我是那种倾向于这样做的生物”。与无模型强化学习中习惯独立于目标导向策略选择而被习得不同，主动推断中的习惯是通过反复追求目标导向策略而形成的，例如把这些策略的结果缓存下来。

在主动推断中，目标导向机制与习惯机制可以合作，而不只是竞争。这是因为，对策略的先验信念同时依赖于一个习惯项，即策略的先验价值，以及一个审议项，即期望自由能。主动推断的层级化扩展甚至暗示，反应式机制与目标导向机制可以被组织成一个层级，而不是并行通路（Pezzulo, Rigoli, and Friston 2015）。

最后还值得指出的是，主动推断与强化学习在如何理解行为及其成因上也有微妙差异。强化学习起源于行为主义理论，认为行为是由强化所中介的试错学习结果。主动推断则假定，行为是推断的结果。这就引向下一小节。

\subsection{规划即推断}

就像知觉问题可以被表述为推断问题一样，控制问题也可以被表述为一种近似贝叶斯推断问题（Todorov 2008）。与此一致，在主动推断中，规划被看作一种推断过程，即对生成模型中的控制状态序列进行推断。

这一想法与其他若干路径密切相关，包括“控制即推断”（Rawlik et al. 2013; Levine 2018）、“规划即推断”（Attias 2003; Botvinick and Toussaint 2012），以及风险敏感控制与 KL 控制（Kappen et al. 2012）。在这些方法中，规划通过推断关于动作或动作序列的后验分布来进行；其依据是一个动态生成模型，该模型编码状态、动作与未来期望状态之间的概率联结。通过在生成模型中对未来奖励（Pezzulo and Rigoli 2011; Solway and Botvinick 2012）或最优未来轨迹（Levine 2018）进行条件化，就可以推断出最佳动作或计划。例如，可以把模型中的未来期望状态夹定，也就是固定其取值，然后推断最可能把当前状态带到未来期望状态的动作序列。

主动推断、规划即推断以及其他相关方案，都采用一种前瞻式控制：它们从对未来待观测状态的显式表征出发，而不是像最优控制与强化学习那样更多地从一组刺激--反应规则或策略出发。不过，控制即推断与规划即推断的具体实现至少沿着三个维度有所不同：使用什么形式的推断，例如抽样还是变分推断；推断的对象是什么，例如关于动作还是动作序列的后验分布；以及推断的目标是什么，例如最大化某种最优性条件的边缘似然，还是最大化获得奖励的概率。

主动推断在这三个维度上都采取了独特立场。第一，它使用一种可扩展的近似方案，即变分推断，来解决规划即推断中出现的困难计算问题。第二，它支持基于模型的规划，也就是对控制状态后验的推断，而控制状态对应的是动作序列或策略，而非单个动作。\footnote{在离散时间主动推断中，策略通常就是未来一段时间内控制状态的序列。} 第三，为了推断动作序列，主动推断考虑的是期望自由能这一泛函；它在数学上囊括了其他广泛使用的规划即推断方案，例如 KL 控制，并且能够处理含糊不清的情形（Friston, Rigoli et al. 2015）。

\section{行为与有限理性}

在主动推断中，行为会自动混合多种成分：审议性、持续性与习惯性（Parr 2020）。设想一个人正在从家附近走向一家商店。如果她能够预测自己行动的后果，例如向左转还是向右转会发生什么，那么她就能形成一个到达商店的好计划。行为中的这一审议性面向由期望自由能提供；当一个人以实现偏好观测的方式行动时，例如出现在商店里，期望自由能就会被最小化。需要注意的是，期望自由能还包含降低不确定性的驱动力，这一点也会体现在审议过程中。例如，如果此人不确定哪条路最好，她可能会先走到一个更合适的观察位置，从那里更容易看清通往商店的路，即便这意味着绕远一些。总之，她的计划会获得认识性可供性。

如果这个人较少有能力进行审议，例如因为她分心了，那么她可能在到达商店后还继续往前走。这种行为中的持续性面向由变分自由能提供；当一个人收集到与当前信念相容的观测时，包括与当前行动进程相关的信念，变分自由能就会被最小化。她所收集到的感觉与本体感觉观测会为“正在行走”这一解释提供证据，因此在缺少审议时便可能导致行为持续。

最后，当这个人不太能进行审议时，她还可能不假思索地选择那条平常回家的路线。这种习惯性成分由策略的先验价值提供。它会把较高概率赋给“回家”这一计划，因为这是她过去多次观察自己实际执行过的计划；在没有被审议过程压过的情况下，这一先验便可能占据主导。

需要注意的是，审议性、持续性与习惯性并非互斥，而是在主动推断中共存并可组合。换言之，主体完全可以推断出：在当前情境中，习惯是最可能的行动路线。这一点不同于那些“双系统理论”，后者假定我们被两个彼此分离的系统所驱动，一个理性，一个直觉（Kahneman 2017）。行为中这三种成分的混合比例，似乎取决于情境条件，例如经验多少，以及个体能投入到审议过程中的认知资源多少，因为审议往往伴随着较高的复杂性成本。\footnote{这里的复杂性成本，可以理解为为了让信念更精确而付出的信息处理代价。}

认知资源对决策的影响，通常在“有限理性”这一框架下被研究（Simon 1990）。其核心思想是：理想的理性主体应当总是充分考虑自身行动的后果，而有限理性的主体则必须平衡计算的成本、努力与时效性，例如为审议最佳计划而承担的信息处理成本（Todorov 2009; Gershman et al. 2015）。

\subsection{有限理性的自由能理论}

有限理性已经被表述为 Helmholtz 自由能最小化问题，而这一热力学构造与主动推断中所使用的变分自由能概念有严格关联；详见 Gottwald and Braun（2020）。“有限理性的自由能理论”用自由能的两个组成部分，即能量与熵，来表述在信息处理能力受限时进行动作选择所面临的权衡，见第 2 章。前者代表一个选择的期望价值，也就是准确性项；后者代表审议成本，也就是复杂性项。审议之所以有成本，是因为为了在决策前让信念更精确，就需要降低其熵，也就是降低其不确定性或复杂性（Ortega and Braun 2013; Zénon et al. 2019）。直观地说，更精确的后验信念会带来更准确的选择，也可能蕴含更高效用；但由于提高信念精度本身是有代价的，因此有限理性的决策者必须通过最小化自由能来寻求折中。在主动推断中，同样的权衡自然会出现，从而生成有限理性的行为形式。有限理性的观念也呼应了主动推断的一个决定性特征，即使用证据，也就是边缘似然的变分下界。总之，主动推断提供了一种关于有限理性与最优性的模型：在给定问题上，最佳解不是单纯把某个单一指标推到极致，而是来自准确性与复杂性之间的折中。这些目标源于一种比经典目标，如经济学常讨论的效用最大化，更丰富的规范性要求，即自由能最小化。

\section{效价、情绪与动机}

主动推断把负自由能看作适应度与生物体实现目标能力的一种度量。虽然主动推断认为生物体会通过行动来最小化自由能，但这并不意味着它们必须显式计算自由能本身。一般来说，只要沿着自由能的梯度行动即可。打个比方，我们不需要知道自己当前的海拔高度，也能通过顺着坡向上走而抵达山顶。不过，也有人提出，生物体可能会建模自身自由能如何随时间变化。支持这一假设的人认为，这将有助于刻画效价、情绪与动机等现象。

按照这一观点，情绪效价，也就是情绪的正向或负向性质，可以被理解为自由能随时间变化的速率，也就是自由能关于时间的一阶导数（Joffily and Coricelli 2013）。具体而言，当一个生物体体验到自由能随时间上升时，它可能会对当下情境赋予负效价；而当它体验到自由能随时间下降时，则可能赋予正效价。若再把这一思路扩展到自由能的长期动态，乃至其二阶时间导数，便可能刻画更复杂的情绪状态，例如从低效价阶段转入高效价阶段时的宽慰，或从高效价阶段跌入低效价阶段时的失望。监控自由能动态及其引发的情绪状态，可能有助于系统根据长期环境统计来调整行为策略或学习速率。

当然，假设还存在第二个生成模型，专门负责监控第一个生成模型的自由能，这看起来似乎跨得太远了。不过，这些思想还可以有另一种解释方式。一个有趣的形式化思路是，去思考究竟是什么导致了自由能的快速变化。由于自由能本身是关于信念的泛函，因此自由能的快速变化必定来自快速的信念更新。而决定这种更新速度的关键因素是精度；在预测编码动力学中，精度相当于一个时间常数。有意思的是，这又把我们带回自由能的高阶导数，因为精度等于自由能景观曲率的负值，也就是二阶导数的负号。不过，这又引出了一个问题：为什么我们应把精度与效价联系起来？答案在于，精度与模糊性成反比。某件事越精确，其解释就越不含糊。选择一条能够最小化期望自由能的行动路线，也就意味着最小化模糊性，从而最大化精度。这样一来，我们就直接把自由能的高阶导数、自由能变化速率与受动机驱动的行为联系起来了。

关于自由能增长或下降的预期，也可能发挥动机功能并激励行为。在主动推断中，一个关于自由能变化的替代性预期量，就是关于策略的信念精度，这再次凸显了这一二阶统计量的重要性。举例而言，高精度信念表明主体已经找到了一条好策略，也就是一条可以被高度自信地预期为最小化自由能的策略。有意思的是，关于策略的信念精度已经被联系到多巴胺信号（FitzGerald, Dolan, and Friston 2015）。从这一角度看，那些提高策略信念精度的刺激会触发多巴胺爆发，而这可能标示了它们的激励显著性（Berridge 2007）。这一视角也许能帮助解释：为什么对目标或奖励即将实现的预期，会与注意增加（Anderson et al. 2011）和动机增强（Berridge and Kringelbach 2011）之间存在神经生理联系。

\section{内稳态、异稳态与内感受加工}

一个生物体的生成模型不仅关乎外部世界，而且也关乎，也许更重要地关乎，其内部环境。关于身体内部的生成模型，也就是内感受图式，承担双重角色：一方面解释内感受，也就是身体感觉，是如何产生的；另一方面确保生理参数得到正确调节（Iodice et al. 2019），例如体温或血糖水平。控制论传统，前文第 10.6.2 节已略有涉及，认为生命体的一个核心目标是维持内稳态（Cannon 1929），即让生理参数保持在可生存范围内，例如体温不能高得离谱；而内稳态只有通过成功控制环境才能实现（Ashby 1952）。

在主动推断中，这种内稳态调节可以通过把生理参数的可存活范围设定为关于内感受观测的先验来实现。有趣的是，内稳态调节可以通过多层嵌套方式完成。最简单的调节回路，是当某些参数已经超出，或预计将超出范围时，启动自主神经反射，例如血管扩张，以应对过高体温。这类自主控制可以被构造为内感受推断，也就是一种作用于内感受流而非本体感觉流的主动推断过程，后者更对应朝向外部的行动（Seth et al. 2012; Seth and Friston 2016; Allen et al. 2019）。为此，大脑可以使用一个预测内感受与生理流的生成模型，并通过触发自主神经反射来修正内感受预测误差，例如“体温出奇地高”。这类似于系统如何通过激活运动反射来修正本体感觉预测误差，并引导朝向外部的行动。

主动推断并不止步于简单的自主回路。它还可以用越来越复杂的方式来修正同一个内感受预测误差，例如体温过高（Pezzulo, Rigoli, and Friston 2015）。它可以采用预测性的异稳态策略（Sterling 2012; Barrett and Simmons 2015; Corcoran et al. 2020），超越单纯内稳态，在内感受预测误差被真正触发之前，就以异稳态方式预先控制生理状态，例如在过热之前先去找阴凉处。另一种预测性策略，是在预计生理设定点将发生偏离之前，预先动员资源，例如在长跑之前提高心输出量，以应对即将增加的氧气需求。这要求系统动态修改关于内感受观测的先验，从而超越静态内稳态（Tschantz et al. 2021）。最终，预测性大脑还可以发展出更复杂的目标导向策略，例如去海滩时提前带上冰水，以更丰富、更高效的方式满足同一个要求，也就是控制体温。

生物学调节与内感受调节，对情感与情绪加工可能至关重要（Barrett 2017）。在情境化互动中，大脑的生成模型不断预测的不只是“接下来会发生什么”，还包括“这些事件对内感受与异稳态会带来什么后果”。在感知外部对象与事件时被激发的内感受流，会为这些对象与事件赋予情感维度，提示它们对生物体的异稳态与生存究竟是好还是坏，从而使它们变得“有意义”。如果这一看法是正确的，那么内感受与异稳态加工的障碍，就可能引发情绪失调以及多种精神病理状态（Pezzulo 2013; Barrett et al. 2016; Maisto, Barca et al. 2019; Pezzulo, Maisto et al. 2019）。

与内感受推断相伴而生的，还有一个新兴方向，即情绪推断。在这种主动推断应用中，情绪被视为生成模型的一部分：它们只是大脑用来在深层生成模型中部署精度的另一类构念或假设。从信念更新的角度看，这意味着焦虑不过是对贝叶斯信念“我很焦虑”的一种承诺，而这一信念最能解释当前占优的感觉与内感受线索。从行动的角度看，随之而来的内感受预测会增强或减弱各种精度，也就是隐性行动，或者驱动自主神经反应，也就是显性行动。这种状态看起来很像唤醒，从而反过来确认“我很焦虑”这一假设。通常，情绪推断涉及一种领域一般性的信念更新，它会同时整合来自内感受与外感受感觉流的信息；这也解释了为何情绪、内感受与注意在健康状态（Seth and Friston 2016; Smith, Lane et al. 2019; Smith, Parr, and Friston 2019）与疾病状态（Peters et al. 2017; J. E. Clark et al. 2018）中关系如此紧密。

\section{注意、显著性与认识动态}

考虑到本章中我们已经无数次提到精度与期望自由能，如果不专门给注意与显著性留出一点篇幅，那就失职了。这两个概念贯穿整个心理学史，且经历过无数次重定义与分类。有时，人们用它们来指涉突触增益控制机制（Hillyard et al. 1998），即优先选择某种感觉模态，或某个模态内的一部分通道；有时，它们又指我们如何通过显性或隐性行动来调整定向，以获取更多世界信息（Rizzolatti et al. 1987; Sheliga et al. 1994, 1995）。

尽管“注意”这一术语的多重含义使该研究领域带有某种认识上的吸引力，但消解由此带来的歧义同样有价值。形式化心理学的一个好处就在于，我们不必被这些歧义困住。我们可以在操作上把注意定义为与某些感觉输入相关联的精度。这一定义与增益控制的概念恰好契合，因为那些被我们推断为更精确的感觉，会比那些被推断为不精确的感觉，对信念更新产生更大影响。这种对应关系的构念效度已经在若干心理学范式中得到展示，包括著名的 Posner 范式（Feldman and Friston 2010）。更具体地说，对视觉空间中被赋予更高精度位置上的刺激所做出的反应，会比对其他位置刺激的反应更快。

这就使“显著性”这一术语也需要一个类似的形式定义。通常，在主动推断中，我们把显著性与期望信息增益，也就是认识价值，联系起来；它是期望自由能的一个组成部分。直观上说，我们越预期某个对象或位置会提供更多信息，它就越显著。不过，这一定义是相对于动作或策略来说的，而注意则是感觉输入信念的属性。这与把显著性理解为显性或隐性定向是一致的。我们在第 7 章中已经看到，还可以把期望信息增益进一步拆分为显著性与新奇性：前者关乎推断的潜力，后者关乎学习的潜力。一个能够表达注意与显著性或新奇性区别的类比，是科学实验的设计与分析。注意，是在我们已经测得的数据中，挑选质量最高的数据，并用它们来支持假设检验的过程；显著性，则是设计下一次实验，以确保能获得质量最高的数据。

我们在这里停留这一问题，并不是要再给文献贡献一种新的注意现象分类，而是要强调：一旦承诺于一种形式化心理学，就会带来一个重要优势。在主动推断之下，别人如何定义注意，或者如何定义任何其他构念，都不再重要，因为我们可以直接诉诸相应的数学构造，从而避免混淆。最后还有一点值得考虑：上述定义还提供了一个简单解释，说明为什么注意与显著性常常被混为一谈。高精度数据意味着最小模糊性。因此，这些数据理应得到注意，而获取这些数据的行动也会显得高度显著（Parr and Friston 2019a）。

\section{规则学习、因果推断与快速泛化}

人类与其他动物擅长进行复杂因果推断，学习抽象概念与对象之间的因果关系，并从有限经验中实现泛化。相比之下，当前机器学习范式往往需要大量样本才能取得类似表现。这种差异提示我们，现有机器学习方法虽然在模式识别方面高度成熟，但可能并没有完整捕捉人类学习与思考的方式（Lake et al. 2017）。

主动推断的学习范式建立在生成模型的发展之上，而生成模型用来刻画动作、事件与观测之间的因果关系。本书中，我们讨论的任务相对简单，例如第 7 章中的 T 形迷宫例子，只需要不太复杂的生成模型即可。相反，要理解并推理复杂情境，就需要深层生成模型，能够把握环境中的潜在结构，例如那些允许我们跨越表面不相似情境进行泛化的隐藏规律（Tervo et al. 2016; Friston, Lin et al. 2017）。

一个由隐藏规则支配复杂社会互动的简单例子，是交通路口。设想一个天真的观察者站在繁忙十字路口，试图预测或解释“在什么情况下行人过马路、车辆通过”。他可以积累关于事件共现的统计，例如一辆红车停下时，一个高个子男人过街；或一位老妇人停下时，一辆大车驶过，但其中大多数统计最终都没什么用。经过一段时间，他也许会发现某些重复出现的统计模式，例如所有车辆在路上某一点停下后不久，行人就会过街。如果任务仅仅是预测行人何时准备过街，那么在机器学习设定中，这种统计规律也许已经足够；但它并不意味着对情境有了真正理解。事实上，它甚至可能导向错误结论，例如以为是“车辆停下”导致了“行人移动”。这类错误在不诉诸因果模型的机器学习应用中很典型，因为它们无法区分是下雨导致草地变湿，还是草地变湿导致下雨（Pearl and Mackenzie 2018）。

另一方面，推断出正确的隐藏规则，例如交通灯规则，则能更深入地理解该情境的因果结构，例如是交通灯导致车辆停下，也导致行人通行。隐藏规则不仅带来更好的预测能力，也使推断更为简约，因为它可以抽离掉大多数感觉细节，例如车辆的颜色。反过来，这又允许系统把理解泛化到其他情境，例如不同的十字路口或不同城市，尽管那些地方的大多数感觉细节都显著不同。当然，面对某些城市的路口，例如罗马，也许只看交通灯还不够。最后，关于交通灯规则的学习，还可能使个体在新情境中的学习更高效，也就是发展出心理学所谓的“学习定势”，或机器学习里的“学会学习”能力（Harlow 1949）。当面对一个交通灯失灵的路口时，先前习得的规则虽然不再能直接使用，但个体仍可能预期：这里还有另一个相似的隐藏规则在起作用，这种预期会帮助他理解交通警察到底在做什么。

这一简单例子说明，学习关于环境潜在结构的丰富生成模型，也就是结构学习，能够支持复杂形式的因果推理与泛化。如何扩展生成模型，使其足以处理这类复杂情境，是计算建模与认知科学中的一个持续目标（Tenenbaum et al. 2006; Kemp and Tenenbaum 2008）。有趣的是，这里存在一种张力：当前机器学习的发展趋势常常强调“越大越好”，而主动推断的统计学方法则强调平衡模型的准确性与复杂性，并偏好更简单的模型。模型约简，也就是修剪不必要的参数，并不仅仅是为了避免浪费资源，它也是学习隐藏规则的一种有效方式，包括在睡眠这类离线阶段中进行（Friston, Lin et al. 2017）；这种过程也许会体现在静息态活动中（Pezzulo, Zorzi, and Corbetta 2020）。

\section{主动推断与其他领域：开放方向}

本书主要关注那些处理生存与适应等生物学问题的主动推断模型。但主动推断同样可以应用于许多其他领域。在最后这一节里，我们简要讨论其中两个领域：社会与文化动力学，以及机器学习与机器人学。前者要求我们思考多个主动推断主体彼此互动的方式，以及这种互动所带来的涌现效应；后者则要求我们理解，如何在不违背该理论基本假设的前提下，为主动推断赋予更高效的学习与推断机制，从而将其扩展到更复杂的问题上。这两者都是极具吸引力的开放研究方向。

\subsection{社会与文化动力学}

人类认知中许多有趣方面，涉及的不是个体主义式的知觉、决策与行动，而是社会与文化动力学（Veissière et al. 2020）。按定义，社会动力学需要多个主动推断生物体彼此互动，这些互动可以是物理性的，例如联合行动或团队运动，也可以是更抽象的，例如选举或社交网络。即使只是相同生物体之间的简单“相互主动推断”演示，也已经产生了有趣的涌现现象，例如能够抵抗离散的简单生命形式的自组织，参与形态发生过程以获得并恢复身体形态的能力，以及相互协调的预测与轮流互动（Friston 2013; Friston and Frith 2015a; Friston, Levin et al. 2015）。其他模拟则关注，生物体如何把自身认知延伸到物质人工物之上，并塑造自己的认知生态位（Bruineberg et al. 2018）。

这些模拟当然只捕捉到了我们社会与文化动力学复杂性的一小部分，但它们已经表明：主动推断具有从“个体科学”扩展为“社会科学”的潜力，也表明认知并不止于头骨之内（Nave et al. 2020）。

\subsection{机器学习与机器人学}

本书讨论的生成建模与变分推断方法，也在机器学习与机器人学中被广泛使用。在这些领域里，重点通常放在如何学习连接主义生成模型上，而不是像本书这样，重点放在如何使用这些模型进行主动推断。这一点很有意义，因为机器学习方法有潜力扩展本书中所讨论的生成模型与问题的复杂度；当然，它们也可能诉诸与主动推断十分不同的过程理论。

虽然我们不可能在这里综述机器学习中庞大的生成建模文献，但仍可简要提及几类最流行的模型，它们也孕育出许多变体。两个早期的连接主义生成模型，Helmholtz 机与 Boltzmann 机（Ackley et al. 1985; Dayan et al. 1995），为如何以无监督方式学习神经网络内部表征提供了范例。Helmholtz 机尤其接近主动推断的变分路线，因为它使用分离的识别网络与生成网络来推断隐变量上的分布，并从这些分布中采样以获得虚构数据。这些早期方法在实践上的成功有限，但后来，多层堆叠的受限 Boltzmann 机使系统能够学习多层内部表征，并成为无监督深度神经网络早期成功的代表之一（Hinton 2007）。

近年的两个连接主义生成模型例子，变分自编码器（variational autoencoders, VAEs）（Kingma and Welling 2014）与生成对抗网络（generative adversarial networks, GANs）（Goodfellow et al. 2014），如今已被广泛用于图像和视频的识别与生成等机器学习应用。VAE 展示了变分方法在生成网络学习中的一种优雅应用。它们的学习目标，即证据下界（evidence lower bound, ELBO），在数学上与变分自由能等价。这个目标既要求对数据做出准确描述，也就是最大化准确性；又偏好那些与先验差异不太大的内部表征，也就是最小化复杂性。后者起到所谓正则化器的作用，有助于泛化并避免过拟合。

GAN 采取的是另一条路径：它由一个生成网络和一个判别网络组成，两者在学习过程中持续竞争。判别网络学习区分生成网络生成的样本数据中哪些是真实数据、哪些是虚构数据；而生成网络则努力生成足以欺骗判别网络，也就是会被其误判的虚构数据。两者之间的竞赛迫使生成网络不断提高自身生成能力，产出高保真虚构数据；这一能力已被广泛用于生成逼真图像等任务。

上述生成模型以及其他模型，也都可以用于控制任务。例如，Ha and Eck（2017）使用一种序列到序列的 VAE 来学习预测铅笔笔画。通过从 VAE 的内部表征中采样，模型可以构造新的笔画式图画。生成建模方法同样被用于控制机器人运动。其中有些方法明确使用主动推断（Pio-Lopez et al. 2016; Sancaktar et al. 2020; Ciria et al. 2021）或与其非常接近的思想，只是处在连接主义设定中（Ahmadi and Tani 2019; Tani and White 2020）。

这一领域的主要挑战之一在于，机器人运动通常是高维的，因此需要学习复杂的生成模型。主动推断及其相关方法的一个有趣之处在于，最需要学习的，是动作与下一时刻感觉反馈，例如视觉与本体感觉反馈之间的前向映射。这种前向映射可以通过多种方式获得：自主探索、示范学习，甚至与人直接互动。例如，教师，也就是实验者，可以抓着机器人的手，沿着通往目标的轨迹引导其移动，从而为高效目标导向动作的习得提供支架（Yamashita and Tani 2008）。能够以多种方式学习生成模型，大大拓展了机器人最终可掌握技能的范围。反过来，利用主动推断发展更先进的神经机器人，不仅对技术发展重要，对理论发展也同样重要。因为主动推断中的一些关键特征，例如适应性的主体--环境互动、认知功能的整合以及具身性的重要性，在机器人场景中都能得到天然检验。

\section{总结}

本书一开始提出的问题是：我们是否可能从第一性原理出发理解大脑与行为？随后，我们把主动推断介绍为一个有资格回应这一挑战的候选理论。我们希望，读者此时已经被说服：对最初问题的回答是肯定的。在本章中，我们考察了主动推断为有感行为提供的统一视角，以及这一理论对熟悉的心理学构念，如知觉、行动选择与情绪，意味着什么。这也让我们有机会重新回顾全书引入的概念，并提醒自己：未来研究仍有许多引人入胜的问题处于开放状态。我们希望，本书能成为主动推断相关著作的有益补充，一方面补充其哲学讨论（Hohwy 2013; Clark 2015），另一方面也补充其物理学讨论（Friston 2019a）。

现在，我们来到了这段旅程的终点。我们的目标，一直是为那些有兴趣使用这些方法的人提供一个入门，无论是在概念层面还是形式层面。不过必须强调的是，主动推断并不是一种只靠阅读理论就能真正掌握的东西。我们鼓励所有喜欢这本书的读者，进一步在实践中追求它。在理论神经生物学中，一些重要的成人礼包括：试着写下一个生成模型；亲身体验模拟表现失常时的挫败；以及当意料之外的事情发生时，学会从自身先验信念被违背之中获得教训。无论你是否选择在计算层面继续推进这种实践，我们都希望，你会在日常生活中反思自己所进行的主动推断。它也许表现为，你不由自主地把目光转向视野边缘某个不确定的对象，以消解不确定性；也许表现为，你选择去一家自己喜欢的餐馆吃饭，以满足先验的味觉偏好；也许表现为，当淋浴水太热时，你会把温度调低，好让世界重新符合你关于其应然样貌的模型。归根结底，我们相信，你会以某种形式持续地实践主动推断。
